\documentclass[10pt]{article}
\usepackage[margin=0.8in]{geometry}
\usepackage{graphicx,booktabs}
\usepackage[T1]{fontenc}
\usepackage{lmodern}
\title{N12: Independent-host replay of the frozen held-out factor-base evaluation}
\author{Reproducibility packet}
\date{27 October 2026}
\begin{document}
\maketitle
\noindent\textbf{Observed result.} The frozen \texttt{fb1-any-j32} policy and matched source-negation rho each returned 420/420 source-verified scalars on seven held-out models. Every target identity and complete operation-count record matches the first host. The seven raw JSONL files are byte identical. At $S/M=0.5$, $I/M=10$, and logical price $0.1$, the ratio of summed candidate cost to summed rho cost is \textbf{2.479905}. The aggregate grid has 0 winning, 36 losing, and 0 tied cells. Individual target ratios fall below one on 19/420 targets.

\begin{center}
\begin{tabular}{rrrrr}
\toprule
Bits & Verified & Ratio of sums & Target wins & Grid wins/losses \\
\midrule
16 & 60/60 & 1.606758 & 0 & 0/36 \\
20 & 60/60 & 1.839211 & 2 & 0/36 \\
24 & 60/60 & 2.334273 & 2 & 0/36 \\
28 & 60/60 & 2.314935 & 2 & 0/36 \\
32 & 60/60 & 2.461524 & 5 & 0/36 \\
36 & 60/60 & 2.247559 & 3 & 0/36 \\
40 & 60/60 & 2.549056 & 5 & 0/36 \\
\midrule
All & 420/420 & 2.479905 & 19 & 0/36 \\
\bottomrule
\end{tabular}
\end{center}

\begin{center}\includegraphics[width=0.96\linewidth]{ratio-by-size.png}\end{center}
\noindent\small Blue points are target-level ratios; orange diamonds are ratios of summed costs by bit size. The vertical scale is logarithmic. The dashed line marks equal cost.\normalsize

\paragraph{Replay boundary.} The original run used EC2 instance \texttt{i-09774109223fde89e}; this replay used EKS node instance \texttt{i-0df0361be2629437d}. Both used the same frozen selector, source workloads, native policy, and price grid. The replay's 11 input members passed a native SHA-256 manifest check before computation. A native validator rebuilt each source target and checked both recovered scalars, phase totals, identity, operation counts, and all 36 aggregate price cells.

\paragraph{Attempt accounting.} Jobs v1 and v2 exited before target execution because the pinned image lacks \texttt{hostname} and \texttt{tar}, respectively. Their exit receipts are retained. Job v3 used per-file S3 transfer and completed 7/7 slices, 60 targets each, with exit code zero. Its 2-CPU, 4-GiB, 3,600-second, no-retry boundary was unchanged.

\paragraph{Executable provenance.} A read-only download of the frozen
\texttt{heldout} from its original S3 input prefix matches the pod's
SHA-256 receipt, \texttt{58a70ac6...e2e0e86e}. The checked-in source files
match their frozen SHA-256 values. An offline, locked release build with the
same rustc 1.99.1 through the ordinary checkout path yielded a different
same-size binary. Building the same source through the original recorded
source path yielded a byte-identical
executable (\texttt{cmp} exit zero). Both build commands and full hashes are
in \texttt{BINARY\_REBUILD.json}. This recheck ran no target trial and did
not change any cost ratio.

\paragraph{Cost and limits.} Costs are cold single-target source-addition equivalents, including construction, transfer, collection, elimination, and source verification. The 36-cell grid uses $S/M\in\{0,0.5,1\}$, $I/M\in\{1,10,100,1000\}$, and logical price $\{0,0.1,1\}$. Wall-time receipts are diagnostic and are not substituted for operation costs. This packet tests the frozen seven-model, 420-target cohort; it does not estimate a population confidence interval.
\end{document}
